\section{Desenvolvendo com Containers}

\begin{frame}
  \frametitle{Iniciando com Docker Compose}
  Para iniciar um projeto existente com Docker:

  \vspace{0.5em}
  \texttt{git clone https://github.com/docker/getting-started-todo-app}

  \vspace{0.3em}
  \texttt{docker compose watch}

  \vspace{0.5em}
  Ao inicializar, vários containers são criados, cada um com uma função específica.
\end{frame}

\begin{frame}
  \frametitle{Containers da Aplicação}
  \begin{itemize}
    \item \textbf{React frontend}: servidor de desenvolvimento React com Vite
    \item \textbf{Node backend}: API para retirar, criar e deletar itens de tarefas
    \item \textbf{MySQL database}: banco de dados da aplicação para armazenar a lista de itens
    \item \textbf{phpMyAdmin}: interface web para o banco em \texttt{http://db.localhost}
    \item \textbf{Traefik proxy}: direciona os \textit{requests} para o serviço correto, tudo acessível pela porta 80
  \end{itemize}
\end{frame}

\begin{frame}
  \frametitle{Fazendo Mudanças no Serviço}
  É possível alterar aspectos da aplicação diretamente nos arquivos do projeto:

  \vspace{0.5em}
  \begin{itemize}
    \item \textbf{Greeting}: alterável em \texttt{backend/src/routes/getGreeting.js}
    \item \textbf{Placeholder text}: alterável em \texttt{client/src/components/AddNewItemForm.jsx}
    \item \textbf{Cor de fundo}: alterável em \texttt{client/src/index.css}
  \end{itemize}
\end{frame}

\begin{frame}
  \frametitle{Vantagens do Ambiente Containerizado}
  \begin{block}{Zero configuração prévia}
    Ambiente de desenvolvimento completo sem instalar Node, MySQL ou qualquer dependência, apenas Docker e um editor de código.
  \end{block}

  \vspace{0.5em}
  \begin{itemize}
    \item Arquivos locais sincronizados automaticamente com os containers
    \item Processos monitoram e respondem a alterações em tempo real
  \end{itemize}
\end{frame}